% !PRINTED 476
\weberpage{476}

As a simple example we mention here only the \gesperrt{group of cyclic
permutations, briefly the cyclic group}, which one obtains when, in any
arrangement of digits $(1, 2, 3 \ldots n)$, one replaces each digit by
the following one and the last by the first, and when one repeats this
permutation until one comes back to the original arrangement. The
degree of the cyclic group is equal to the number of digits, e.\,g.
with three digits:
\[
\begin{pmatrix} 1, & 2, & 3 \\ 2, & 3, & 1 \end{pmatrix}, \quad
\begin{pmatrix} 1, & 2, & 3 \\ 3, & 1, & 2 \end{pmatrix}, \quad
\begin{pmatrix} 1, & 2, & 3 \\ 1, & 2, & 3 \end{pmatrix}.
\]

There are special groups, which will occupy us still further later
on, that consist only of such permutations in whose composition,
besides the laws mentioned above, the commutative law also holds, for
which therefore always $\pi_1 \pi_2 = \pi_2 \pi_1$, and in consequence
of which the analogy with multiplication is a much more complete one.
Such groups, to which the cyclic groups belong, are called
\gesperrt{commutative} or \gesperrt{Abelian groups}. When all the
elements of a permutation group $Q_1$ are contained in a group $Q$,
then $Q_1$ is called a \gesperrt{divisor of $Q$, and $Q$ divisible by
$Q_1$}.

\begin{center}
\S.\,149.

\gesperrt{Galois Group.}
\end{center}

If, in the roots $\alpha, \alpha_1, \alpha_2 \ldots \alpha_{m-1}$ of
an equation $F(x) = 0$, as they are expressed by \S.\,148, (3) as
quantities of the normal field $\Omega(\varrho)$ of degree $\mu$, we
carry out one of the $\mu$ substitutions $\sigma$ of the normal field,
then a certain permutation takes place among these roots, as has
already been shown above. To every substitution $\sigma$ there
corresponds therefore a certain permutation $\pi$ of the $m$ digits
$0, 1, 2 \ldots m-1$. \gesperrt{Two different substitutions $\sigma_h, \sigma_k$
lead also to two different permutations $\pi_h, \pi_k$.}

For by \S.\,145, $\varrho$ is a rational function of $\alpha,
\alpha_1 \ldots \alpha_{m-1}$; we use the following notation.
\[
\varrho = \Phi(\alpha, \alpha_1 \ldots \alpha_{m-1})
\]
If we here set
\[
\alpha_s = \chi_s(\varrho)
\]
it follows that

\weq{1}{\varrho = \Phi[\chi(\varrho), \chi_1(\varrho) \ldots \chi_{m-1}(\varrho)],}
%CONT>>
